\problemname{Cherimoyas}

Farah loves the exotic cherimoya fruit. Since it comes from a
faraway land, it is only sold one day per year in Sweden! Farah has of course
made sure to buy some cherimoyas on this very day.

The cherimoyas are at different stages of ripeness. Some have just become ripe on the same day while
others will become ripe to eat later.

More precisely, each cherimoya fruit is ripe to eat for a total of 3 days. We
say that the fruit becomes ripe to eat on that day. Before that it cannot be eaten and after the
three days it must be thrown away.

Farah wants to get as much as possible out of the cherimoya season. She wants to
maximize her enjoyment, which is calculated as follows: For a given day she receives 10
enjoyment points for the first cherimoya, then 9 for the second, 8 for the
third, and so on. She never manages to eat more than 10 cherimoyas in a day.

Write a program that, given how many cherimoyas become ripe to eat each day,
determines how many enjoyment points Farah will get during the year's
cherimoya season.

\section*{Input}
The first line of input contains the integer $N$ ($1 \leq N \leq 15$), the number
of days on which cherimoyas can become ripe.

For each of the following $N$ days you get an integer $C_i$ ($0 \leq C_i \leq 30$),
meaning that $C_i$ cherimoyas become ripe on that day.
It is thus during $N+2$ days that eating cherimoyas is relevant.

\section*{Output}
Output a single integer: the maximum number of enjoyment points Farah can
get with the best eating strategy.

\section*{Scoring}
Your solution will be tested on a set of test groups.
To earn points for a group, you must pass all test cases in that group.

\noindent
\begin{tabular}{| l | l | p{12cm} |}
  \hline
  \textbf{Group} & \textbf{Points} & \textbf{Constraints} \\ \hline
  $1$    & $20$       & $N \leq 1$ \\ \hline
  $2$    & $20$       & $N \leq 2$ \\ \hline
  $3$    & $30$       & $N \leq 5$ \\ \hline
  $4$    & $30$       & No additional constraints. \\ \hline
\end{tabular}

\section*{Explanation of Sample 1}
An optimal strategy is the following:
\begin{enumerate}
    \item Of the 18 cherimoyas that ripen on day 1, eat 6 on each of days 1, 2, and 3.
    This gives $10+9+8+7+6+5=45$ enjoyment points on each of those days.
    \item Of the 2 cherimoyas that ripen on day 2, eat one on day 4 and one on day 5.
    This gives $10$ enjoyment points on each of those days.
    \item In total this gives $45 \cdot 3+10 \cdot 2=155$ enjoyment points. There is no strategy
    that gives more enjoyment points.
\end{enumerate}
